\section{2018 Geometry \& Topology}

\subsection{Individual Exam}

\begin{exercise}
\label{geometry:2018:1}
Let $M$ and $N$ be smooth, connected, orientable $n$-manifolds for $n \geq 3$, and let $M \# N$ denote their connected sum.
\begin{enumerate}
    \item[(a)] Compute the fundamental group of $M \# N$ in terms of that of $M$ and of $N$ (you may assume that the basepoint is on the boundary sphere along which we glue $M$ and $N$).
    \item[(b)] Compute the homology groups of $M \# N$.
    \item[(c)] For part (a), what changes if $n = 2$? Use this to describe the fundamental groups of orientable surfaces.
\end{enumerate}
\end{exercise}

\begin{solution}
We use the standard convention that the connected sum is formed from closed connected orientable $n$-manifolds by deleting an open ball from each and gluing the resulting boundary spheres.

\textbf{(a) Fundamental group for $n\geq 3$.}
Write
\[
M^\circ=M\setminus \operatorname{int}D^n,\qquad
N^\circ=N\setminus \operatorname{int}D^n.
\]
Then
\[
M\# N=M^\circ\cup_{S^{n-1}}N^\circ.
\]
For $n\geq 3$, the gluing sphere $S^{n-1}$ is simply connected. Also removing a ball from an $n$-manifold with $n\geq 3$ does not change the fundamental group, so
\[
\pi_1(M^\circ)\cong \pi_1(M),\qquad
\pi_1(N^\circ)\cong \pi_1(N).
\]
By Seifert--van Kampen,
\[
\pi_1(M\# N)\cong \pi_1(M)*\pi_1(N).
\]

\textbf{(b) Homology.}
For reduced homology, apply Mayer--Vietoris to
\[
M\# N=M^\circ\cup N^\circ,\qquad M^\circ\cap N^\circ\simeq S^{n-1}.
\]
For $1\leq k\leq n-2$, the group $H_k(S^{n-1})$ vanishes, so exactness gives
\[
H_k(M\# N)\cong H_k(M^\circ)\oplus H_k(N^\circ)
\cong H_k(M)\oplus H_k(N).
\]
The top dimension is
\[
H_n(M\# N)\cong \mathbb{Z}
\]
because the connected sum is again closed, connected, and orientable.

In degree $n-1$, exactness gives
\[
0\to H_n(M\# N)\to H_{n-1}(S^{n-1})
\to H_{n-1}(M^\circ)\oplus H_{n-1}(N^\circ)
\to H_{n-1}(M\# N)\to 0.
\]
The first nonzero map sends the fundamental class of $M\#N$ to the boundary sphere with opposite orientations on the two punctured pieces. Equivalently, it accounts for the single relation that glues the two deleted balls. The result is
\[
H_{n-1}(M\#N)\cong H_{n-1}(M)\oplus H_{n-1}(N).
\]
Thus, for closed connected orientable manifolds,
\[
H_0(M\# N)\cong \mathbb{Z},\qquad
H_n(M\# N)\cong \mathbb{Z},
\]
and for $1\leq k\leq n-1$,
\[
H_k(M\#N)\cong H_k(M)\oplus H_k(N).
\]

\textbf{(c) The case $n=2$.}
Now the gluing sphere is $S^1$, which is not simply connected. Seifert--van Kampen no longer gives a free product without relations; the boundary loop of the punctured surface must be identified on the two sides.

For an orientable surface $\Sigma_g$ of genus $g$,
\[
\Sigma_g=T^2\#\cdots\#T^2
\]
with $g$ torus summands. Starting from the once-punctured torus presentation and imposing the final boundary relation gives
\[
\pi_1(\Sigma_g)
\cong
\left\langle
a_1,b_1,\dots,a_g,b_g
\ \middle|\
\prod_{i=1}^g [a_i,b_i]=1
\right\rangle.
\]
So connected sum of surfaces adds generators but also keeps track of the boundary commutator relation.
\end{solution}

\begin{exercise}
\label{geometry:2018:2}
Determine all of the possible degrees of maps $S^2 \to S^1 \times S^1$.
\end{exercise}

\begin{solution}
The degree of a map between closed oriented manifolds of the same dimension is defined by the induced map on top homology.

Let
\[
f:S^2\to S^1\times S^1=T^2.
\]
Since
\[
H_2(S^2;\mathbb{Z})\cong \mathbb{Z},
\qquad
H_2(T^2;\mathbb{Z})\cong \mathbb{Z},
\]
the degree is the integer $d$ satisfying
\[
f_*[S^2]=d[T^2].
\]

But $f$ must have degree zero. One way to see this is to use cohomology. Let $\alpha,\beta\in H^1(T^2;\mathbb{Z})$ be the two standard degree-one generators. Then
\[
H^2(T^2;\mathbb{Z})=\mathbb{Z}\langle \alpha\smile\beta\rangle.
\]
Because
\[
H^1(S^2;\mathbb{Z})=0,
\]
we have
\[
f^*\alpha=f^*\beta=0.
\]
Therefore
\[
f^*(\alpha\smile\beta)=f^*\alpha\smile f^*\beta=0.
\]
On top cohomology, pulling back the fundamental class of $T^2$ multiplies by the degree, so this says $d=0$.

The constant map realizes degree $0$. Hence the only possible degree is
\[
0.
\]
\end{solution}

\begin{exercise}
\label{geometry:2018:3}
Classify all vector bundles over the circle $S^1$ up to isomorphism.
\end{exercise}

\begin{solution}
We classify real vector bundles of fixed rank $r$ over $S^1$.

Cut the circle into two intervals whose intersection has two contractible components. Since every vector bundle over an interval is trivial, a rank $r$ bundle over $S^1$ is determined by a clutching function on the overlap. Equivalently, after identifying $S^1$ with $[0,1]/(0\sim 1)$, it is determined by one transition matrix
\[
A\in GL(r,\mathbb{R}),
\]
gluing
\[
(0,v)\sim (1,Av).
\]
Changing trivializations conjugates and homotopes this matrix, so the isomorphism class depends only on the connected component of $GL(r,\mathbb{R})$ containing $A$.

The group $GL(r,\mathbb{R})$ has two connected components, distinguished by the sign of the determinant. Thus there are exactly two rank $r$ real vector bundles over $S^1$:
\[
\det A>0
\quad\text{and}\quad
\det A<0.
\]
The first is the trivial bundle
\[
S^1\times \mathbb{R}^r.
\]
The second is the non-orientable bundle, represented by
\[
[0,1]\times \mathbb{R}^r /(0,v)\sim (1,Bv),
\qquad
B=\operatorname{diag}(-1,1,\dots,1).
\]
Equivalently, it is the direct sum of the Möbius line bundle and a trivial rank $r-1$ bundle.

For rank $0$ there is only the zero bundle. For complex vector bundles, the corresponding clutching group is $GL(r,\mathbb{C})$, which is connected, so every complex vector bundle over $S^1$ is trivial.
\end{solution}

\begin{exercise}
\label{geometry:2018:4}
Suppose $C$ is a regular curve in the unit sphere $S^2$. For any point $W \in S^2$, there exists the only oriented great circle $S_W$ (determined by the right hand rule) in $S^2$ such that $W$ is the pole of $S_W$. Denote by $n(W)$ the number of points at which the oriented great circle $S_W$ and $C$ intersect. Prove the Crofton formula
\[
\iint_{S^2} n(W) dW = 4L,
\]
where $dW$ and $L$ are the area element of $S^2$ and the length of $C$, respectively.
\end{exercise}

\begin{solution}
It is enough to prove the formula first for a single smooth embedded arc and then add over a partition of the curve; transverse intersection failures occur only on a measure-zero set of poles $W$.

Parametrize $C$ by arc length:
\[
\gamma:[0,L]\to S^2,\qquad |\gamma'(s)|=1.
\]
The great circle with pole $W$ is
\[
S_W=\{X\in S^2:\langle X,W\rangle=0\}.
\]
Thus $S_W$ meets $C$ exactly at the zeros of
\[
F(s,W)=\langle \gamma(s),W\rangle.
\]

For fixed $s$, the condition $F(s,W)=0$ says that $W$ lies on the great circle
\[
\gamma(s)^\perp\cap S^2.
\]
For almost every $W$, the intersections are transverse. For such a $W$, the one-dimensional delta-function form of the coarea formula gives
\[
n(W)=\int_0^L
\delta(\langle \gamma(s),W\rangle)
|\langle \gamma'(s),W\rangle|\,ds.
\]
Integrating this identity over $S^2$ and using Fubini gives
\[
\int_{S^2}n(W)\,dW
=
\int_0^L
\left(
\int_{S^2}
\delta(\langle \gamma(s),W\rangle)
|\langle \gamma'(s),W\rangle|\,dW
\right)ds.
\]
Now fix $s$. The delta factor restricts $W$ to the level set
\[
\langle \gamma(s),W\rangle=0.
\]
On this level set the spherical gradient of $W\mapsto \langle \gamma(s),W\rangle$ has norm $1$, so the delta integral becomes an ordinary line integral over the corresponding great circle.
Choose coordinates at each $s$ so that
\[
\gamma(s)=e_3,\qquad \gamma'(s)=e_1.
\]
Then the pole circle is
\[
W=(\cos\theta,\sin\theta,0),
\qquad 0\leq \theta<2\pi,
\]
and
\[
|\langle \gamma'(s),W\rangle|=|\cos\theta|.
\]
The inner integral is
\[
\int_{S^2}\delta(\langle \gamma(s),W\rangle)
|\langle \gamma'(s),W\rangle|\,dW
=\int_0^{2\pi}|\cos\theta|\,d\theta=4.
\]
Therefore each unit length of curve contributes $4$ to the incidence count. Integrating over $s$ gives
\[
\int_{S^2} n(W)\,dW
=\int_0^L 4\,ds
=4L.
\]
This is the Crofton formula.
\end{solution}

\begin{exercise}
\label{geometry:2018:5}
Let $M$ be an $n$-dimensional closed submanifold in the Euclidean space $\mathbb{R}^{n+p}$. Prove the following inequality
\[
\int_{M} H^n dV \geq \mathrm{vol}(S^n),
\]
where $H$ and $dV$ are the mean curvature (i.e., norm of the mean curvature vector) and the volume element of $M$, and $S^n$ is the standard unit sphere of dimension $n$.
\end{exercise}

\begin{solution}
This is the Chern--Lashof--Reilly inequality. The argument below records the standard proof chain. The integral-geometric Chern--Lashof estimate is a deep input and is not proved from first principles here.

\textbf{Step 1: Use height functions.}
For each unit vector $u\in S^{n+p-1}$, consider the height function
\[
h_u:M\to \mathbb{R},\qquad h_u(x)=\langle x,u\rangle.
\]
Since $M$ is compact, $h_u$ attains a maximum and a minimum. For almost every $u$, $h_u$ is Morse. Hence it has at least one critical point of index $0$ and at least one critical point of index $n$.

The critical point condition is exactly
\[
u\in \nu_xM,
\]
where $\nu_xM$ is the normal space of $M$ at $x$.

\textbf{Step 2: Apply the normal Gauss map.}
Let
\[
UN(M)=\{(x,u):x\in M,\ u\in \nu_xM,\ |u|=1\}
\]
be the unit normal bundle, and let
\[
G:UN(M)\to S^{n+p-1},\qquad G(x,u)=u.
\]
The preceding paragraph says that almost every $u\in S^{n+p-1}$ has at least two preimages under $G$, counted with suitable Morse index information. The Chern--Lashof integral-geometric argument refines this count and gives
\[
\int_M K^*(x)\,dV\geq \operatorname{vol}(S^n),
\]
where $K^*(x)$ is the average over the unit normal sphere of the absolute Gauss--Kronecker curvature in that normal direction.

\textbf{Step 3: Compare the averaged Gauss--Kronecker curvature with $H^n$.}
Let $A_\xi$ be the shape operator in a unit normal direction $\xi$. The normal curvature determinant in direction $\xi$ is
\[
|\det A_\xi|.
\]
The arithmetic-geometric mean inequality, applied to the principal normal curvatures and then averaged over normal directions, gives the pointwise estimate used in Reilly's proof:
\[
K^*(x)\leq H(x)^n,
\]
where $H(x)$ is the norm of the mean curvature vector under the convention in the problem.

Combining the two inequalities gives
\[
\int_M H^n\,dV
\geq
\int_M K^*\,dV
\geq
\operatorname{vol}(S^n).
\]
This proves the desired inequality modulo the quoted Chern--Lashof integral-geometric estimate. Equality is attained by the round unit sphere in an affine $(n+1)$-plane.
\end{solution}

\begin{exercise}
\label{geometry:2018:6}
Let $M$ be an even dimensional compact and oriented Riemannian manifold with positive sectional curvature. Show that $M$ is simply connected.
\end{exercise}

\begin{solution}
This is the even-dimensional part of Synge's theorem.

\textbf{Synge's theorem.}
If $M$ is compact, orientable, even-dimensional, and has positive sectional curvature, then $M$ is simply connected.

\textbf{Proof.}
Assume for contradiction that $\pi_1(M)$ is nontrivial. Let $\alpha$ be a nontrivial free homotopy class of closed curves. Since $M$ is compact, the length functional attains its minimum in this class, so there is a closed geodesic
\[
\gamma:[0,\ell]\to M
\]
representing $\alpha$ and minimizing length in its free homotopy class.

Let $P$ be parallel transport around $\gamma$ on the normal space
\[
\dot\gamma(0)^\perp.
\]
Because $\dim M$ is even, this normal space has odd dimension. Since $M$ is orientable and $\gamma$ is closed, the parallel transport map $P$ preserves orientation. An orientation-preserving orthogonal transformation of an odd-dimensional Euclidean space has eigenvalue $1$. Thus there exists a nonzero vector
\[
v\in \dot\gamma(0)^\perp
\]
with $P(v)=v$.

Parallel translate $v$ along $\gamma$ to obtain a parallel normal field $V(t)$ satisfying
\[
V(0)=V(\ell).
\]
Therefore $V$ is an admissible periodic variation field through closed curves in the same free homotopy class.

The second variation of length for a closed geodesic in a periodic normal variation is
\[
I(V,V)=\int_0^\ell
\left(|\nabla_{\dot\gamma}V|^2
-\langle R(V,\dot\gamma)\dot\gamma,V\rangle
\right)dt.
\]
Since $V$ is parallel, the first term vanishes. Positive sectional curvature gives
\[
\langle R(V,\dot\gamma)\dot\gamma,V\rangle>0
\]
for all $t$, so
\[
I(V,V)<0.
\]
This means that $\gamma$ can be varied through nearby closed curves in the same free homotopy class to obtain shorter curves, contradicting the choice of $\gamma$ as a length minimizer.

Hence no nontrivial free homotopy class exists, so $\pi_1(M)=0$.
\end{solution}

\subsection{Team Exam}

\begin{exercise}
\label{geometry:2018:7}
Let $X$ be $(S^{2} \times S^{2}) \cup_{S^{2}} D^{3}$, where we attach the 3-disk via the map
\[
S^{2} \longrightarrow S^{2} \vee S^{2}
\]
which crushes a great circle connecting the north and south poles. Compute the homology groups of $X$.
\end{exercise}

\begin{solution}
\textbf{Prerequisites:}
\begin{itemize}
    \item \textbf{CW Complex Structure}: A way to build spaces by gluing "cells" (disks) of various dimensions together.
    \item \textbf{Relative Homology \& Long Exact Sequence}: A sequence that relates the homology of a space $X$, a subspace $M$, and the relative homology $H_k(X, M)$.
    \item \textbf{Attaching Cells}: When we attach a $k$-cell to a space $M$ via a map $S^{k-1} \to M$, the relative homology $H_k(X, M)$ is $\mathbb{Z}$, generated by the new cell, and the boundary map $\partial$ gives the degree of the attaching map.
\end{itemize}

\textbf{Intuition:}

The space $M = S^2 \times S^2$ can be visualized by its cellular structure: one $0$-cell, two $2$-cells (one for each $S^2$ factor), and one $4$-cell. The 2-skeleton is thus $S^2 \vee S^2$.

We are given a space $X$ formed by taking $M$ and attaching a solid 3-dimensional disk $D^3$ to it. The boundary of $D^3$ is a 2-sphere $S^2$. The problem states that the attaching map from $S^2$ to $M$ crushes a great circle connecting the north and south poles. 

When you pinch a great circle on a sphere to a point, the sphere becomes two spheres joined at a point, which is exactly $S^2 \vee S^2$. Therefore, the attaching map essentially wraps the boundary of the 3-disk exactly once around each of the two $S^2$ factors of $M$. 

By adding a 3-cell, we are introducing a new potential 3-cycle, but we are also introducing a relation in the 2-cycles. The boundary of this new 3-cell is the sum of the two $S^2$ generators. This relation will "kill" one of the 2-dimensional homology generators.

\textbf{Steps:}

\textbf{Step 1: Identify the homology of the base space $M = S^2 \times S^2$.}

Using the K\"unneth formula or cellular homology, the homology groups of $M$ are:
$H_0(M) = \mathbb{Z}$, $H_1(M) = 0$, $H_2(M) = \mathbb{Z} \oplus \mathbb{Z}$, $H_3(M) = 0$, and $H_4(M) = \mathbb{Z}$.

Let $a$ and $b$ be the generators of $H_2(M)$ corresponding to the two $S^2$ factors.

\textbf{Step 2: Understand the boundary map of the new 3-cell.}

The space $X$ is formed by $X = M \cup_\phi e^3$, where the attaching map is $\phi: S^2 \to M$. 

The relative homology groups $H_k(X, M)$ are zero for all $k \neq 3$, and $H_3(X, M) \cong \mathbb{Z}$, generated by the new 3-cell $e^3$.

The boundary map $\partial: H_3(X, M) \to H_2(M)$ in the long exact sequence of the pair $(X, M)$ is simply the map induced by the attaching map $\phi_*: H_2(S^2) \to H_2(M)$.

Because $\phi$ crushes a great circle, it maps the fundamental class of $S^2$ to the sum of the fundamental classes of the two spheres in $S^2 \vee S^2$. Thus, $\partial(1) = \phi_*([S^2]) = a + b$.

\textbf{Step 3: Apply the long exact sequence of the pair $(X, M)$.}

The long exact sequence is:
$\dots \to H_k(M) \to H_k(X) \to H_k(X, M) \xrightarrow{\partial} H_{k-1}(M) \to \dots$

For $k=0, 1$:
Since $H_0(X, M) = 0$ and $H_1(X, M) = 0$, $H_0(X) \cong H_0(M) \cong \mathbb{Z}$ and $H_1(X) \cong H_1(M) = 0$.

For $k=2$:
We have the sequence $H_3(X, M) \xrightarrow{\partial} H_2(M) \to H_2(X) \to H_2(X, M) = 0$.

Substituting the known groups: $\mathbb{Z} \xrightarrow{\partial} \mathbb{Z} \oplus \mathbb{Z} \to H_2(X) \to 0$.

The image of $\partial$ is the subgroup generated by $(1, 1)$ since $\partial(1) = a + b$. 

By exactness, $H_2(X) \cong (\mathbb{Z} \oplus \mathbb{Z}) / \text{im}(\partial) \cong (\mathbb{Z} \oplus \mathbb{Z}) / \langle (1, 1) \rangle \cong \mathbb{Z}$.

For $k=3$:
We have the sequence $0 = H_4(X, M) \to H_3(M) \to H_3(X) \to H_3(X, M) \xrightarrow{\partial} H_2(M)$.

Substituting the known groups: $0 \to 0 \to H_3(X) \to \mathbb{Z} \xrightarrow{\partial} \mathbb{Z} \oplus \mathbb{Z}$.

Since $\partial(1) = (1, 1)$, the map $\partial$ is injective. Therefore, its kernel is $0$.

By exactness, $H_3(X)$ must be isomorphic to the kernel of $\partial$, so $H_3(X) = 0$.

For $k=4$:
We have $0 = H_5(X, M) \to H_4(M) \to H_4(X) \to H_4(X, M) = 0$.

Thus, $H_4(X) \cong H_4(M) \cong \mathbb{Z}$.

\textbf{Conclusion:}
The homology groups of $X$ are $H_0(X) = \mathbb{Z}$, $H_1(X) = 0$, $H_2(X) = \mathbb{Z}$, $H_3(X) = 0$, and $H_4(X) = \mathbb{Z}$.
\end{solution}

\begin{exercise}
\label{geometry:2018:8}
(a) Let $A$ be a single circle in $\mathbb{R}^{3}$. Compute the fundamental group $\pi_{1}(\mathbb{R}^{3} - A)$.

(b) Let $A$ and $B$ be disjoint circles in $\mathbb{R}^{3}$, supported in the upper and lower half space, respectively. Compute $\pi_{1}(\mathbb{R}^{3} - (A \cup B))$.
\end{exercise}

\begin{solution}
\textbf{Pedagogical Note:} This problem explores the topology of complements of knots and links. Even for unknots, the complement has interesting structure.

\textbf{Geometric Intuition:}
Imagine $\mathbb{R}^3$ as $S^3$ minus a point at infinity. A circle $A$ in $\mathbb{R}^3$ can be viewed as an unknotted circle in $S^3$. The complement of a circle in $S^3$ is another circle (they are linked in the sense of the Hopf fibration). When we remove a point (infinity) from this complement circle, it doesn't change the homotopy type from $S^1$ unless the point was on the circle itself.

\textbf{Formal Proof:}
(a) Consider $S^3 = \mathbb{R}^3 \cup \{\infty\}$. The complement of a circle $A$ in $S^3$ is homotopy equivalent to $S^1$ (think of $S^3$ as two solid tori $S^1 \times D^2 \cup D^2 \times S^1$ glued along their boundary $S^1 \times S^1$; if $A$ is the core of one torus, the complement is the other torus, which retracts to its core $S^1$).
Then $\mathbb{R}^3 \setminus A = S^3 \setminus (A \cup \{\infty\})$. This is homotopy equivalent to $S^1 \vee S^2$ (removing a point from a 3-manifold is like adding a 2-sphere at that point).
However, for fundamental group purposes, $\pi_1(\mathbb{R}^3 \setminus A) \cong \pi_1(S^1 \vee S^2) \cong \pi_1(S^1) * \pi_1(S^2) \cong \mathbb{Z} * \{e\} \cong \mathbb{Z}$.

(b) Let $U = \{z > - \epsilon\} \setminus A$ and $V = \{z < \epsilon\} \setminus B$. Then $\mathbb{R}^3 \setminus (A \cup B) = U \cup V$. The intersection $U \cap V \cong \mathbb{R}^2 \times (-\epsilon, \epsilon)$ is contractible, so $\pi_1(U \cap V) = 0$. 
By the Seifert-van Kampen theorem:
\[
\pi_1(\mathbb{R}^3 \setminus (A \cup B)) \cong \pi_1(U) *_{\pi_1(U \cap V)} \pi_1(V) \cong \pi_1(U) * \pi_1(V) \cong \mathbb{Z} * \mathbb{Z} \cong F_2
\]
where $F_2$ is the free group on two generators. Intuitively, each circle provides one independent loop that cannot be shrunk to a point.
\end{solution}

\begin{exercise}
\label{geometry:2018:9}
Consider the differential 1-form $\omega = x\, dy - y\, dx + dz$ in $\mathbb{R}^{3}$ with coordinates $(x, y, z)$. Prove that $f\omega$ is not closed for any nowhere zero function $f : \mathbb{R}^{3} \to \mathbb{R}$.
\end{exercise}

\begin{solution}
\textbf{Prerequisites:}
\begin{itemize}
    \item \textbf{Differential Forms}: Objects that can be integrated over manifolds. A 1-form acts on vectors, and a 3-form acts on volumes.
    \item \textbf{Exterior Derivative ($d$)}: A generalization of the gradient, curl, and divergence. For a function $f$, $df$ is its gradient form. For a 1-form $\omega$, $d\omega$ is a 2-form. A form $\omega$ is \textit{closed} if $d\omega = 0$.
    \item \textbf{Wedge Product ($\wedge$)}: A way to multiply differential forms. It is antisymmetric for 1-forms, meaning $dx \wedge dy = -dy \wedge dx$ and $dx \wedge dx = 0$.
    \item \textbf{Product Rule for Forms}: $d(f\omega) = df \wedge \omega + f d\omega$.
\end{itemize}

\textbf{Intuition:}

The problem asks us to show that no matter how we scale the 1-form $\omega$ by a non-zero function $f$, the result $f\omega$ can never be "closed" (meaning its exterior derivative is zero).

If a 1-form could be scaled to be closed, it would locally look like the gradient of some function $g$, meaning it would be orthogonal to the level surfaces of $g$. However, the given 1-form $\omega = x\, dy - y\, dx + dz$ is a classic example of a \textit{contact form}. 

A contact form describes a "twisted" plane field that constantly rotates as you move through space. Because it twists so much, there are no 2D surfaces everywhere tangent to it. The algebraic test for this maximum twistiness is computing $\omega \wedge d\omega$. If this 3-form is non-zero everywhere, the form cannot be scaled to be closed.

\textbf{Steps:}

\textbf{Step 1: Compute the exterior derivative of $\omega$.}

The form is $\omega = x\, dy - y\, dx + dz$. We take the exterior derivative $d$ of each term.

$d(x\, dy) = dx \wedge dy + x\, d(dy) = dx \wedge dy + 0 = dx \wedge dy$.

$d(-y\, dx) = -dy \wedge dx - y\, d(dx) = -dy \wedge dx = dx \wedge dy$ (using the antisymmetric property).

$d(dz) = 0$.

Adding these together, we get $d\omega = dx \wedge dy + dx \wedge dy = 2\, dx \wedge dy$.

\textbf{Step 2: Assume for contradiction that $f\omega$ is closed.}

Suppose there exists a nowhere zero function $f$ such that $d(f\omega) = 0$.

By the product rule for exterior derivatives, $d(f\omega) = df \wedge \omega + f d\omega = 0$.

This means $df \wedge \omega = -f d\omega$.

\textbf{Step 3: Show that this leads to a contradiction using the wedge product.}

Let's take the wedge product of both sides of the equation with $\omega$ from the left:
$\omega \wedge (df \wedge \omega) = \omega \wedge (-f d\omega)$.

Because the wedge product is associative and $\omega \wedge \omega = 0$ for any 1-form (since any 1-form wedged with itself is zero), the left side becomes:
$\omega \wedge df \wedge \omega = -\omega \wedge \omega \wedge df = 0$.

So we must have $0 = -f (\omega \wedge d\omega)$. 

Since we assumed $f$ is nowhere zero, this implies we must have $\omega \wedge d\omega = 0$ everywhere.

\textbf{Step 4: Compute $\omega \wedge d\omega$ to find the contradiction.}

We have $\omega = x\, dy - y\, dx + dz$ and $d\omega = 2\, dx \wedge dy$.

Let's compute their wedge product:
$\omega \wedge d\omega = (x\, dy - y\, dx + dz) \wedge (2\, dx \wedge dy)$.

Distributing the wedge product:
$= 2x(dy \wedge dx \wedge dy) - 2y(dx \wedge dx \wedge dy) + 2(dz \wedge dx \wedge dy)$.

Notice that $dy \wedge dy = 0$ and $dx \wedge dx = 0$, so the first two terms vanish.

We are left with $\omega \wedge d\omega = 2(dz \wedge dx \wedge dy)$.

Rearranging the differentials (which requires two swaps: $dz$ past $dx$, then $dz$ past $dy$, so the sign stays the same):
$\omega \wedge d\omega = 2\, dx \wedge dy \wedge dz$.

This is twice the standard volume form in $\mathbb{R}^3$, which is non-zero everywhere. 

This contradicts the requirement that $\omega \wedge d\omega = 0$. Therefore, no such function $f$ can exist.
\end{solution}

\begin{exercise}
\label{geometry:2018:10}
Show that
\[
Q^{n} := \left\{(x^{1}, \dots, x^{n+1}) \in \mathbb{R}^{n+1}; \sum_{i=1}^{n+1} (x^{i})^{4} = 1 \right\}
\]
is a differentiable manifold.
\end{exercise}

\begin{solution}
\textbf{Geometric Intuition:}
$Q^n$ is a "squircle-like" generalization to $n$ dimensions. Just as $x^2+y^2=1$ is a circle, $x^4+y^4=1$ is a smooth closed curve that looks more like a square with rounded corners.

\textbf{Formal Proof:}
Define $g: \mathbb{R}^{n+1} \to \mathbb{R}$ by $g(x) = \sum (x^i)^4$. $Q^n = g^{-1}(1)$.
$1$ is a regular value if $dg \neq 0$ for all $x \in g^{-1}(1)$.
\[
dg = \sum_{i=1}^{n+1} 4(x^i)^3 dx^i
\]
The differential $dg$ vanishes only if $4(x^i)^3 = 0$ for all $i$, i.e., at the origin $(0, \dots, 0)$.
However, $g(0, \dots, 0) = 0 \neq 1$. So $dg$ is never zero on $Q^n$.
By the Regular Value Theorem, $Q^n$ is a smooth manifold of dimension $n$.
\end{solution}

\begin{exercise}
\label{geometry:2018:11}
Let $M$ be a closed surface in $\mathbb{R}^{3}$. Prove that
\[
\int_{M} |K|\, d\sigma \geq 4\pi(1 + g),
\]
where $K$, $g$ and $d\sigma$ is the Gaussian curvature, the genus and the area element of $M$, respectively.
\end{exercise}

\begin{solution}
\textbf{Prerequisites:}
\begin{itemize}
    \item \textbf{Gaussian Curvature ($K$)}: A measure of how much a surface bends. A sphere has positive curvature, a flat plane has zero curvature, and a saddle shape has negative curvature.
    \item \textbf{Euler Characteristic ($\chi$)}: A topological number for a surface. For a closed surface with $g$ "holes" (genus), it is given by $\chi = 2 - 2g$.
    \item \textbf{Gauss-Bonnet Theorem}: A beautiful theorem connecting geometry to topology. It states that the total Gaussian curvature of a closed surface $M$ is exactly $2\pi$ times its Euler characteristic: $\int_M K \, d\sigma = 2\pi \chi(M) = 4\pi(1 - g)$.
    \item \textbf{Gauss Map}: A map that takes every point on a surface to its unit normal vector, placing that vector on the unit sphere $S^2$. The integral of positive curvature measures the area of the image of the Gauss map on $S^2$.
\end{itemize}

\textbf{Intuition:}

The Gauss-Bonnet theorem tells us that the "net" curvature of a surface (positive curvature minus the absolute value of the negative curvature) is fixed by its topology (its number of holes). 

However, we are asked to find a bound on the \textit{absolute} total curvature, which treats both "bulging out" (positive) and "saddle-like" (negative) bending as positive contributions. 

If you put a surface inside $\mathbb{R}^3$, it has to curve back on itself to close up. Think of a potato: no matter how you dent it, there are parts of it facing outward in every possible 3D direction. The normal vectors pointing outward must cover the entire unit sphere at least once. This implies that the total positive curvature must be at least the area of a sphere, which is $4\pi$. 

We can then use simple algebra to relate the absolute curvature to the total positive curvature and the Gauss-Bonnet constant.

\textbf{Steps:}

\textbf{Step 1: Separate the curvature into positive and negative parts.}

Let $M_+$ be the regions on the surface where the curvature is positive or zero ($K \geq 0$).
Let $M_-$ be the regions where the curvature is strictly negative ($K < 0$).

By the Gauss-Bonnet theorem, the integral of $K$ over the whole surface $M$ depends only on its genus $g$:
$\int_M K \, d\sigma = \int_{M_+} K \, d\sigma + \int_{M_-} K \, d\sigma = 4\pi(1 - g)$.

\textbf{Step 2: Express the total absolute curvature in terms of these parts.}

The quantity we want to bound is the integral of the absolute value $|K|$. 
Since $|K| = K$ on $M_+$ and $|K| = -K$ on $M_-$, we can write:
$\int_M |K| \, d\sigma = \int_{M_+} K \, d\sigma - \int_{M_-} K \, d\sigma$.

\textbf{Step 3: Eliminate the negative curvature term.}

From the Gauss-Bonnet equation in Step 1, we can solve for the negative part:
$-\int_{M_-} K \, d\sigma = \int_{M_+} K \, d\sigma - 4\pi(1 - g)$.

Substitute this into our absolute curvature equation from Step 2:
$\int_M |K| \, d\sigma = \int_{M_+} K \, d\sigma + \left( \int_{M_+} K \, d\sigma - 4\pi(1 - g) \right)$.

Simplifying this yields:
$\int_M |K| \, d\sigma = 2 \int_{M_+} K \, d\sigma - 4\pi(1 - g) = 2 \int_{M_+} K \, d\sigma - 4\pi + 4\pi g$.

\textbf{Step 4: Bound the positive curvature using the Gauss map.}

For any closed surface $M$ contained in $\mathbb{R}^3$, if you look at it from the outside, there must be at least one point where the surface's normal vector points in any given direction. (Imagine moving a plane inward from infinity from any direction until it touches the surface; the tangent point has that plane's normal).

This means the Gauss map, which takes points on $M$ to their unit normal vectors on the sphere $S^2$, must cover the entire sphere $S^2$ at least once just using points where the surface curves like a dome (positive curvature). 

Therefore, the total positive curvature must be at least the area of the unit sphere:
$\int_{M_+} K \, d\sigma \geq \text{Area}(S^2) = 4\pi$.

\textbf{Step 5: Conclude the final inequality.}

Substitute this bound into our equation from Step 3:
$\int_M |K| \, d\sigma \geq 2(4\pi) - 4\pi + 4\pi g$.

Simplify the expression:
$\int_M |K| \, d\sigma \geq 8\pi - 4\pi + 4\pi g = 4\pi + 4\pi g$.

Factoring out $4\pi$, we get:
$\int_M |K| \, d\sigma \geq 4\pi(1 + g)$.
\end{solution}

\begin{exercise}
\label{geometry:2018:12}
Let $M$ be an $n$-dimensional compact and simply connected Riemannian manifold. If the sectional curvature $K_{M}$ of $M$ satisfies
\[
\frac{1}{4} < K_{M} \leq 1,
\]
then $M$ is homeomorphic to $S^{n}$.
\end{exercise}

\begin{solution}
\textbf{Context:} 
This is the \textit{Quarter-Pinch Sphere Theorem}. The constant $1/4$ is sharp, as $\mathbb{CP}^n$ with the Fubini-Study metric has sectional curvature in the range $[1/4, 1]$. Since $\pi_2(\mathbb{CP}^n) = \mathbb{Z}$, it is not a sphere.

\textbf{Proof:}
The result is precisely the classical topological sphere theorem of Rauch--Berger--Klingenberg:

\[
\text{If }M\text{ is compact, simply connected, and }0<\delta<K_M\leq 1
\text{ with }\delta>\frac14,\text{ then }M\text{ is homeomorphic to }S^n.
\]

Thus the stated assumptions imply the conclusion by taking any constant $\delta$ satisfying
\[
\frac14<\delta<\inf_M K_M.
\]
For completeness, we recall what the theorem proves and where the strict quarter-pinching enters.

\begin{enumerate}
    \item Rauch comparison controls Jacobi fields along geodesics. The upper bound $K_M\leq 1$ gives a model scale, while the lower bound $K_M>\delta>1/4$ gives strict focusing estimates.

    \item Berger's estimates and Klingenberg's injectivity-radius analysis use simple connectedness to rule out the short geodesic-loop obstruction. This is the hard global step; it is not the elementary statement that the injectivity radius is automatically at least $\pi$.

    \item The comparison estimates imply that geodesic spheres and cut-locus structure have the same topological behavior as in the round sphere. One then constructs a twisted-sphere decomposition by two geodesic balls and proves it is homeomorphic to the standard sphere.
\end{enumerate}

Therefore $M$ is homeomorphic to $S^n$ by the Rauch--Berger--Klingenberg topological sphere theorem. This solution leaves that theorem as an explicit deep dependency. The stronger differentiable sphere theorem, proved later by Brendle--Schoen using Ricci flow, upgrades the conclusion to diffeomorphism under the same strict pointwise quarter-pinching.
\end{solution}
